\chapter{فرض طاعة رسول الله ومكانة السنة من القرآن}

\section{طاعة الرسول صلى الله عليه وسلم مطلقة ومستقلة}
الحمد لله والصلاة والسلام على رسول الله، أما بعد؛ فقد عقد الإمام الشافعي رحمه الله تعالى في كتاب «الرسالة» باباً في بيان فرض طاعة رسول الله صلى الله عليه وسلم، سواء أكانت مقرونة بطاعة الله عز وجل أم كانت مذكورة وحدها.

فالله عز وجل فرض طاعة نبيه صلى الله عليه وسلم في كل حال، وسواء كان النبي صلى الله عليه وسلم مبيناً لما جاء عن ربه، أو أتى بحكم مستقل ليس منصوصاً عليه في كتاب الله، فالفرض علينا أن نطيعه. قال تعالى: ﴿وَمَا كَانَ لِمُؤْمِنٍ وَلَا مُؤْمِنَةٍ إِذَا قَضَى اللَّهُ وَرَسُولُهُ أَمْرًا أَن يَكُونَ لَهُمُ الْخِيَرَةُ مِنْ أَمْرِهِمْ ۗ وَمَن يَعْصِ اللَّهَ وَرَسُولَهُ فَقَدْ ضَلَّ ضَلَالًا مُّبِينًا﴾. فلا خيار لأهل الإيمان في طاعة رسول الله صلى الله عليه وسلم متى ما صح عنه الأمر.

\section{طاعة أولي الأمر مقيدة ومستثناة}
قال تعالى: ﴿يَا أَيُّهَا الَّذِينَ آمَنُوا أَطِيعُوا اللَّهَ وَأَطِيعُوا الرَّسُولَ وَأُولِي الْأَمْرِ مِنكُمْ﴾. يُلاحظ أن الله عز وجل كرر لفظ "أطيعوا" مع الرسول صلى الله عليه وسلم للدلالة على استقلال طاعته، بينما لم يكرره مع "أولي الأمر"؛ لأن طاعة أولي الأمر (من الأمراء والعلماء) مقيدة بطاعة الله وطاعة رسوله.

فلا طاعة لمخلوق في معصية الخالق، وقد قال النبي صلى الله عليه وسلم: «إنما الطاعة في المعروف». فطاعة أولي الأمر ليست طاعة مطلقة في كل شيء، بل هي طاعة مستثناة؛ فإذا أمروا بمعصية فلا سمع لهم ولا طاعة في تلك المعصية. ولهم علينا حق السمع والطاعة وعدم الخروج عليهم ما أقاموا شرع الله، وعليهم إقامة الدين وحفظ الثغور وأعراض ودماء المسلمين. 

\section{الرد عند التنازع وحجية القياس}
قال تعالى: ﴿فَإِن تَنَازَعْتُمْ فِي شَيْءٍ فَرُدُّوهُ إِلَى اللَّهِ وَالرَّسُولِ إِن كُنتُمْ تُؤْمِنُونَ بِاللَّهِ وَالْيَوْمِ الْآخِرِ﴾. الاختلاف واقع بين الأمة وولاة أمرها وعلمائها، ولكن المرجعية عند التنازع يجب أن تكون لكتاب الله وسنة رسوله صلى الله عليه وسلم على فهم سلف الأمة، وليس إلى آراء البشر أو القوانين الوضعية والمناهج المنحرفة كالديمقراطية أو العلمانية.

فإن لم نجد نصاً ظاهراً في مسألة التنازع، رددناه قياساً على كتاب الله وسنة رسوله. والقياس نوعان:
\begin{itemize}
    \item \textbf{قياس صحيح:} وهو القياس الذي فيه إعمال للنصوص الشرعية، كإلحاق فرع بأصل في حكم لعلة جامعة بينهما (مثل القياس في تحديد القبلة أو معرفة العدل).
    \item \textbf{قياس باطل فاسد الاعتبار:} وهو القياس المصادم للنصوص، كقياس إبليس عندما أُمر بالسجود لآدم فقال: ﴿أَنَا خَيْرٌ مِّنْهُ﴾، فقاس مع وجود النص الصريح بالسجود فكان قياساً باطلاً.
\end{itemize}

\section{وجوه السنة مع القرآن الكريم}
السنة النبوية مع كتاب الله عز وجل تأتي على وجوه:
\begin{enumerate}
    \item \textbf{نص موافق:} أن تنزل السنة مؤكدة وموافقة لما جاء به نص القرآن الكريم (كفرضية الصلاة والزكاة).
    \item \textbf{بيان للمجمل:} أن يأتي الفرض مجملاً في القرآن، فتبين السنة كيفيته وتفاصيله عن الله عز وجل، مثل مواقيت الصلوات وعدد ركعاتها، وكيفية الزكاة وأنصبتها.
    \item \textbf{تشريع مستقل:} أن يسن النبي صلى الله عليه وسلم ما ليس فيه نص في كتاب الله (كتحريم الجمع بين المرأة وعمتها أو خالتها، وتحريم كل ذي ناب من السباع).
\end{enumerate}
وهذا الوجه الثالث اختلف فيه العلماء في تأصيله: هل شرعه النبي صلى الله عليه وسلم بما أعطاه الله من الحكمة وبما فرض من طاعته؟ أم أن لكل سنة أصلاً في القرآن يعلمه من يعلمه ويجهله من يجهله؟ أم أنها وحي ألقي في روعه صلى الله عليه وسلم؟ 
والخلاصة: أياً كان التوجيه، فإن طاعة رسول الله صلى الله عليه وسلم فرض محتم، وكل ما جاء به وجب قبوله، ولا عذر لأحد في مخالفة أمر عَرَفَ أنه من أمر رسول الله صلى الله عليه وسلم.

\section{الرد على القرآنيين منكري السنة}
إن من يدعي الاكتفاء بالقرآن ليرد سنة النبي صلى الله عليه وسلم هو في الحقيقة مكذب للقرآن ومتبع لهواه؛ لأن القرآن نفسه أمر بطاعة الرسول وحذر من مخالفته. قال الشافعي: استدللنا بآية ﴿وَأَنزَلَ اللَّهُ عَلَيْكَ الْكِتَابَ وَالْحِكْمَةَ﴾ على أن الحكمة هي سنة رسول الله صلى الله عليه وسلم؛ لأنها قُرنت بالكتاب.
فمن قَبِلَ عن الله فرائضه، يجب أن يقبل عن رسول الله سننه؛ لأن الله هو الذي افترض طاعته. ومن فرّق بينهما فقد ردّ على الله جل وعلا، وسنة رسول الله صلى الله عليه وسلم مبيّنة عن الله مراده.

\section{فتاوى فقهية منثورة في ختام الدرس}
أجاب الشيخ حفظه الله في نهاية الدرس عن جملة من الأسئلة الفقهية، نلخصها فيما يلي:
\begin{itemize}
    \item \textbf{قصر وقت الليل والجمع بين الصلوات:} سُئل عمن يعيش في بلاد لا يتجاوز فيها الليل أربع ساعات، ويشق عليه انتظار العشاء لعمله بالنهار. الجواب: المشقة تجلب التيسير، ويجوز له جمع العشاء مع المغرب جمع تقديم إذا شق عليه ذلك، والأولى أن ينظم وقت نومه ليدرك الصلوات في وقتها خروجاً من الخلاف.
    \item \textbf{التوبة من المال الحرام المستثمر:} سُئل عمن كسب مالاً من تجارة الخمر واشترى به داراً، ثم تاب وأراد استثمار الدار. الجواب: كان الأولى التخلص منه بالكلية إن أمكن، وإذا استثمر هذا المال وجب عليه أن يقدر رأس المال الحرام ويجعله ديناً في ذمته، فيتخلص مما يقابله من أرباح بالصدقة بنية التخلص، وله أن ينفق على نفسه وعياله من الحلال المتبقي قدر الحاجة.
    \item \textbf{الجمع للمسجون المستضعف:} سُئل عمن سُجن في بلاد الكفر ويشق عليه أداء كل صلاة في وقتها بسبب التضييق. الجواب: يجوز له الجمع للمشقة، فـ "المشقة تجلب التيسير" والله لا يريد بعباده العسر.
\end{itemize}